\documentclass{article}

\usepackage{amsmath,amssymb}
\usepackage{xurl}
\usepackage[hidelinks]{hyperref}

% Shared commands for notes in docs/paper-gaps/.

\DeclareUnicodeCharacter{2080}{\ifmmode{}_{0}\else\textsubscript{0}\fi}
\DeclareUnicodeCharacter{2081}{\ifmmode{}_{1}\else\textsubscript{1}\fi}
\DeclareUnicodeCharacter{2082}{\ifmmode{}_{2}\else\textsubscript{2}\fi}
\DeclareUnicodeCharacter{2099}{\ifmmode{}_{n}\else\textsubscript{n}\fi}

\newcommand{\Lean}{\textsc{Lean}}
\ExplSyntaxOn
\NewDocumentCommand{\leanid}{m}{
  \ifmmode
    \textsf{\detokenize{#1}}
  \else
    \group_begin:
    \urlstyle{sf}
    \tl_set:Nn \l_tmpa_tl {#1}
    \tl_replace_all:Nnn \l_tmpa_tl {\_} {_}
    \exp_args:NV \path \l_tmpa_tl
    \group_end:
  \fi
}
\ExplSyntaxOff
\providecommand{\tr}{\operatorname{tr}}
\newcommand{\tnleanIssueBase}{https://github.com/LionSR/TNLean/issues/}
\newcommand{\tnleanPrBase}{https://github.com/LionSR/TNLean/pull/}
\newcommand{\tnleanDocBase}{https://sirui-lu.com/TNLean/docs/}
\newcommand{\ghissue}[1]{\href{\tnleanIssueBase#1}{GitHub issue \##1}}
\newcommand{\ghpr}[1]{\href{\tnleanPrBase#1}{PR \##1}}
\newcommand{\doclink}[1]{\href{\tnleanDocBase#1.html}{\path{#1}}}

% Dirac notation and the identity operator, as in blueprint/src/macros/common.tex.
% \providecommand keeps note-local definitions authoritative.
\usepackage{dsfont}
\providecommand{\ket}[1]{|#1\rangle}
\providecommand{\bra}[1]{\langle #1|}
\providecommand{\braket}[2]{\langle #1 | #2 \rangle}
\providecommand{\ketbra}[2]{|#1\rangle\!\langle #2|}
\providecommand{\Id}{\mathds{1}}
\providecommand{\one}{\mathds{1}}
\providecommand{\C}{\mathbb{C}}
\providecommand{\MN}[1]{M_{#1}(\C)}
\providecommand{\MD}{M_D(\C)}

% External referencing for paper sources.  The build helper writes sanitized
% auxiliary files under build/paper-aux/ and excludes duplicated source labels.
\usepackage{xr}

% Import a generated paper auxiliary file with a caller-supplied label prefix.
% The candidate paths support builds from docs/paper-gaps/, the repository root,
% and build/paper-aux/ itself.
\newcommand{\importpaperaux}[2]{%
  \IfFileExists{../../build/paper-aux/#2.aux}{%
    \externaldocument[#1]{../../build/paper-aux/#2}%
  }{%
    \IfFileExists{build/paper-aux/#2.aux}{%
      \externaldocument[#1]{build/paper-aux/#2}%
    }{%
      \IfFileExists{#2.aux}{%
        \externaldocument[#1]{#2}%
      }{}%
    }%
  }%
}

% General external reference macros
\newcommand{\paperref}[2]{\ref{#1-#2}}
\newcommand{\papereqref}[2]{(\ref{#1-#2})}

% CPSV16 source (1606.00608 / MPDO-22-12-17-2.tex) external document and shortcuts
\importpaperaux{cpsv16-}{1606.00608/MPDO-22-12-17-2-tnlean}
\newcommand{\cpsvref}[1]{\paperref{cpsv16}{#1}}
\newcommand{\cpsveqref}[1]{\papereqref{cpsv16}{#1}}

% Verdict marker: \gapnote{<kind>}{<status>}. Kind is the policy
% classification (clarification, local-correction, scope-restriction,
% unfaithful, false-source, open-gap); status is open, wip (elimination
% underway), resolved, or historical. Machine-read by the site index;
% prints a verdict line.
\newcommand{\gapnote}[2]{\par\noindent\textbf{Verdict:} #1 (#2).\par}


\title{The Proportionality Scalar in the MPS Reduction Theorem}
\author{}
\date{2026-08-29}

\begin{document}
\maketitle
\gapnote{false-source}{open}

\section{The printed statement}

Let $A$ be a normal matrix product state tensor and let $B$ be another tensor
with the same physical alphabet.  Proposition~20 of
Ref.~\cite{Molnar2018SemiInjective} assumes that, for some
$\lambda\in\mathbb C$,
\footnote{Source coordinates in arXiv:1706.07329v2:
    \path{cornerproblem.tex:3124--3133}.}
\begin{align}
    V_n(B)
        &=\lambda^n V_n(A)
        \qquad\text{for every positive chain length }n.
    \label{eq:mgsc18_reduction_proportionality_printed_hypothesis}
\end{align}
It concludes that there are rectangular matrices $V,W$ such that
\begin{align}
    VW
        &=1,
    \label{eq:mgsc18_reduction_proportionality_printed_retraction}\\
    VB^{i_1}\cdots B^{i_m}W
        &=A^{i_1}\cdots A^{i_m}
        \qquad\text{for every word }(i_1,\ldots,i_m).
    \label{eq:mgsc18_reduction_proportionality_printed_word_relation}
\end{align}
The source proof subsequently says that $A$ and $B$ generate the same state,
rather than proportional state families, and uses that equality in the
construction of $V,W$ and in the residual-algebra argument
\cite{Molnar2018SemiInjective}.
\footnote{The equality is asserted at
    \path{cornerproblem.tex:3839}, \path{cornerproblem.tex:3865}, and
    \path{cornerproblem.tex:3973--3976}; the proof of Proposition~20 occupies
    \path{cornerproblem.tex:3815--3938}.}

\section{Counterexamples}

The printed implication is false for every nonzero $\lambda\ne1$.  Fix such a
scalar and take one bond coordinate and one physical letter, with
\begin{align}
    A^1
        &=1,
    & B^1
        &=\lambda.
    \label{eq:mgsc18_reduction_proportionality_nonzero_example}
\end{align}
The transfer map of $A$ is the identity on the one-dimensional matrix algebra,
and is therefore primitive.  Moreover,
\begin{align}
    V_n(B)
        &=\lambda^n
    \label{eq:mgsc18_reduction_proportionality_nonzero_left_state}\\
        &=\lambda^n V_n(A),
    \label{eq:mgsc18_reduction_proportionality_nonzero_hypothesis}
\end{align}
so
Eq.~\ref{eq:mgsc18_reduction_proportionality_printed_hypothesis} holds.  But
Eq.~\ref{eq:mgsc18_reduction_proportionality_printed_retraction} would give
\begin{align}
    VB^1W
        &=\lambda VW
    \label{eq:mgsc18_reduction_proportionality_nonzero_reduction_first_step}\\
        &=\lambda,
    \label{eq:mgsc18_reduction_proportionality_nonzero_reduction_value}
\end{align}
whereas
Eq.~\ref{eq:mgsc18_reduction_proportionality_printed_word_relation} requires
$VB^1W=A^1=1$.  Taking $\lambda=2$ gives the simplest nondegenerate numerical
instance.

The zero scalar is additionally degenerate.  Taking $A^1=1$, $B^1=0$, and
$\lambda=0$ satisfies the positive-length hypothesis, but every pair obeying
$VW=1$ has
\begin{align}
    VB^1W
        &=0
    \label{eq:mgsc18_reduction_proportionality_zero_reduction_value}\\
        &\ne A^1.
    \label{eq:mgsc18_reduction_proportionality_zero_contradiction}
\end{align}

\section{Corrected statements}

The proof printed in Ref.~\cite{Molnar2018SemiInjective} establishes the
equality specialization $\lambda=1$.  If $\lambda\ne0$, define
$\widetilde B^i=\lambda^{-1}B^i$.  Then
\begin{align}
    V_n(\widetilde B)
        &=V_n(A).
    \label{eq:mgsc18_reduction_proportionality_rescaled_state_equality}
\end{align}
Applying the equality specialization to $\widetilde B$ gives $VW=1$ and
\begin{align}
    V\widetilde B^{i_1}\cdots\widetilde B^{i_m}W
        &=A^{i_1}\cdots A^{i_m}.
    \label{eq:mgsc18_reduction_proportionality_rescaled_word_relation}
\end{align}
Equivalently, the corrected conclusion for the original tensor is
\begin{align}
    VB^{i_1}\cdots B^{i_m}W
        &=\lambda^m A^{i_1}\cdots A^{i_m}.
    \label{eq:mgsc18_reduction_proportionality_corrected_word_relation}
\end{align}
Thus one may either assume $\lambda=1$ and retain the printed conclusion, or
assume $\lambda\ne0$ and use
Eq.~\ref{eq:mgsc18_reduction_proportionality_corrected_word_relation}.  No
corresponding conclusion follows from the printed hypotheses when
$\lambda=0$.

\section{MPU specialization}

The MPU application in Ref.~\cite{FrancoRubio2025MPUGauging} starts from the
linear representation identity $U_gU_h=U_{gh}$.  Hence the product tensor and
the target tensor generate exactly the same operator at every positive chain
length: this is the case $\lambda=1$.  The fusion equations used there are
therefore unaffected by the defect in the proportional version of
Proposition~20.

The formal development should state the reusable reduction theorem first for
equal positive-length matrix product state families.  A later proportional
version may be derived by the rescaling above, with the power
$\lambda^m$ retained explicitly.  The equality-case implementation and its
selected-reduction residual algebra are tracked by
\href{https://github.com/LionSR/TNLean/issues/7322}{TNLean issue 7322}.

\section{Verdict}

The unscaled conclusion of Proposition~20 is false for every unrestricted
choice of $\lambda\ne1$, not merely in the degenerate zero case.  The source
proof establishes the $\lambda=1$ specialization.  For nonzero $\lambda$, the
corrected general word relation is
Eq.~\ref{eq:mgsc18_reduction_proportionality_corrected_word_relation}; the
zero case requires separate hypotheses.  The MPU use is source-faithful
because it supplies exact equality.

\bibliographystyle{alpha}
\bibliography{references}

\end{document}
